\chapter{LAB3 --- Incident Response: HELIX Systems}

\section{LAB3 --- Incident Response · HELIX Systems}

{\footnotesize\begin{longtable}[]{@{}
  >{\raggedright\arraybackslash}p{0.22\linewidth}
  >{\raggedright\arraybackslash}p{0.72\linewidth}@{}}
\toprule\noalign{}
\endhead
\bottomrule\noalign{}
\endlastfoot
\textbf{Difficulty} & Medium \\
\textbf{Category} & Incident Response · Log Forensics · Blue Team \\
\textbf{Flags} & 3 \\
\textbf{Est.\ time} & 90--120 min \\
\textbf{Key skills} & Log analysis, SSH brute-force forensics, privilege escalation detection, lateral movement \\
\textbf{Nodes} & pfSense · VyOS · Server-Web · Server-DB \\
\end{longtable}}

\href{https://github.com/TheEgea/TFG/tree/main/src/eve-ng/configs/nodes/Lab3}{Lab folder on GitHub} \href{https://github.com/TheEgea/TFG/raw/main/src/eve-ng/topologies/LAB3-IncidentResponse-HELIX.unl}{Download topology (.unl)} \href{../../assets/exercises/lab3-exercise.pdf}{Exercise sheet} \href{../../assets/exercises/lab3-solution.pdf}{Solution guide}

\begin{center}\rule{0.5\linewidth}{0.5pt}\end{center}

\subsection{Scenario}

\textbf{HELIX Systems} is a fictional company with a two-tier internal infrastructure. Unlike the previous labs, the student here plays the \textbf{defender} role: the attack has already happened.

An unknown threat actor brute-forced SSH credentials on the web server, escalated privileges via a misconfigured \texttt{sudoers} entry, established persistence with a backdoor account, and moved laterally to the database server to exfiltrate sensitive client data.

The task is to connect to the affected systems, collect and correlate log evidence, and reconstruct the attacker's complete activity timeline as a SOC analyst.

\begin{center}\rule{0.5\linewidth}{0.5pt}\end{center}

\subsection{Topology}

\begin{figure}[H]
\centering
\includegraphics[width=0.95\textwidth]{lab3-topology.png}
\caption{EVE-NG topology --- Lab 3}
\end{figure}

\begin{lstlisting}
Homelab LAN (192.168.0.0/24)
        |
[pfSense-LAB3]  WAN: 192.168.0.x (DHCP)  $\leftarrow$ student entry (DNAT TCP:22 $\rightarrow$ Server-Web)
        |        LAN: 172.16.2.1/30
[VyOS-LAB3]
  eth0: 172.16.2.2/30     $\leftarrow$ uplink to pfSense
  eth1: 192.168.50.1/24   $\leftarrow$ Net-Servers
  eth2: 192.168.60.1/24   $\leftarrow$ Net-Internal
        |              |
[Server-Web]      [Server-DB]
192.168.50.10     192.168.60.10
Primary target    Lateral movement
auth.log          clients.db
bash history      exfil_marker
\end{lstlisting}

\subsection{Nodes}

{\footnotesize\begin{longtable}[]{@{}
  >{\raggedright\arraybackslash}p{0.20\linewidth}
  >{\raggedright\arraybackslash}p{0.20\linewidth}
  >{\raggedright\arraybackslash}p{0.28\linewidth}
  >{\raggedright\arraybackslash}p{0.22\linewidth}@{}}
\toprule\noalign{}
\begin{minipage}[b]{\linewidth}\raggedright
Node
\end{minipage} & \begin{minipage}[b]{\linewidth}\raggedright
OS
\end{minipage} & \begin{minipage}[b]{\linewidth}\raggedright
IP
\end{minipage} & \begin{minipage}[b]{\linewidth}\raggedright
Role
\end{minipage} \\
\midrule\noalign{}
\endhead
\bottomrule\noalign{}
\endlastfoot
pfSense-LAB3 & pfSense CE & WAN: 192.168.0.x / LAN: 172.16.2.1/30 & Perimeter firewall · DNAT student entry \\
VyOS-LAB3 & VyOS rolling & 172.16.2.2 / 192.168.50.1 / 192.168.60.1 & Router · SNAT \\
Server-Web & Ubuntu Server 24.04 & 192.168.50.10 & Primary forensic target (helix-web) \\
Server-DB & Ubuntu Server 24.04 & 192.168.60.10 & Lateral movement target (helix-db) \\
\end{longtable}}

\subsection{Network segments}

{\footnotesize\begin{longtable}[]{@{}>{\raggedright\arraybackslash}p{0.22\linewidth}
  >{\raggedright\arraybackslash}p{0.22\linewidth}
  >{\raggedright\arraybackslash}p{0.22\linewidth}
  >{\raggedright\arraybackslash}p{0.22\linewidth}@{}}
\toprule\noalign{}
Segment & Subnet & Gateway & Purpose \\
\midrule\noalign{}
\endhead
\bottomrule\noalign{}
\endlastfoot
Net-Link & 172.16.2.0/30 & 172.16.2.1 & pfSense $\leftrightarrow$ VyOS \\
Net-Servers & 192.168.50.0/24 & 192.168.50.1 & Server-Web \\
Net-Internal & 192.168.60.0/24 & 192.168.60.1 & Server-DB \\
Homelab & 192.168.0.0/24 & 192.168.0.1 & WAN · Student entry \\
\end{longtable}}

\begin{center}\rule{0.5\linewidth}{0.5pt}\end{center}

\subsection{Learning objectives}

\begin{enumerate}
\def\labelenumi{\arabic{enumi}.}
\tightlist
\item
  Read and correlate \texttt{auth.log} entries to identify SSH brute-force patterns
\item
  Detect privilege escalation via \texttt{sudo} misconfiguration in \texttt{sudoers}
\item
  Identify persistence mechanisms (backdoor accounts, cron jobs)
\item
  Trace lateral movement between servers using network logs and bash history
\item
  Reconstruct a complete attacker timeline from pre-staged evidence
\end{enumerate}

\begin{center}\rule{0.5\linewidth}{0.5pt}\end{center}

\subsection{Investigation flow}

\begin{lstlisting}
1. SSH into Server-Web via pfSense DNAT (student entry point) --- `ssh analyst@<pfSense-WAN-IP>`
2. Analyse /var/log/auth.log       $\rightarrow$ brute-force source IP + success timestamp
3. Review /home/devops/.bash_history $\rightarrow$ privilege escalation commands
4. Check /etc/sudoers              $\rightarrow$ misconfigured entry (sudo find $\rightarrow$ root)
5. Find /etc/passwd + shadow       $\rightarrow$ backdoor account created post-compromise
6. SSH to Server-DB (192.168.60.10)$\rightarrow$ lateral movement confirmed
7. Locate /opt/helix/data/clients.db $\rightarrow$ exfiltrated database
8. Find /opt/helix/data/.exfil_marker $\rightarrow$ exfiltration evidence
\end{lstlisting}

\begin{center}\rule{0.5\linewidth}{0.5pt}\end{center}

\subsection{Lab startup checklist}

\begin{enumerate}
\def\labelenumi{\arabic{enumi}.}
\item
  Start all nodes from EVE-NG web UI
\item
  Verify pfSense DNAT rule forwards TCP 22 (WAN) $\rightarrow$ Server-Web (192.168.50.10)
\item
  Student entry: \texttt{ssh\ analyst@\textless{}pfSense-WAN-IP\textgreater{}} --- password: \texttt{An@lyst2024}
\item
  Evidence is pre-staged --- no additional setup required

  \begin{itemize}
  \tightlist
  \item
    \textbf{Server-Web entry}: \texttt{analyst} / \texttt{An@lyst2024} (via pfSense DNAT)
  \item
    \textbf{Server-DB}: reachable from Server-Web once lateral movement path is found
  \end{itemize}
\end{enumerate}

\begin{center}\rule{0.5\linewidth}{0.5pt}\end{center}

\subsection{Evidence artefacts}

{\footnotesize\begin{longtable}[]{@{}
  >{\raggedright\arraybackslash}p{0.17\linewidth}
  >{\raggedright\arraybackslash}p{0.54\linewidth}
  >{\raggedright\arraybackslash}p{0.19\linewidth}@{}}
\toprule\noalign{}
\begin{minipage}[b]{\linewidth}\raggedright
Artefact
\end{minipage} & \begin{minipage}[b]{\linewidth}\raggedright
Location
\end{minipage} & \begin{minipage}[b]{\linewidth}\raggedright
Contains
\end{minipage} \\
\midrule\noalign{}
\endhead
\bottomrule\noalign{}
\endlastfoot
\texttt{auth.log} & \path{/var/log/auth.log} (Server-Web) & SSH brute-force + successful login \\
\path{bash_history} & \path{/home/devops/.bash_history} & Post-exploitation commands \\
\texttt{clients.db} & \path{/opt/helix/data/clients.db} (Server-DB) & Exfiltrated client data \\
\path{exfil_marker} & \path{/opt/helix/data/.exfil_marker} (Server-DB) & Exfiltration timestamp \\
\end{longtable}}

\begin{center}\rule{0.5\linewidth}{0.5pt}\end{center}

\subsection{Files}

{\footnotesize\begin{longtable}[]{@{}
  >{\raggedright\arraybackslash}p{0.45\linewidth}
  >{\raggedright\arraybackslash}p{0.45\linewidth}@{}}
\toprule\noalign{}
\begin{minipage}[b]{\linewidth}\raggedright
File
\end{minipage} & \begin{minipage}[b]{\linewidth}\raggedright
Description
\end{minipage} \\
\midrule\noalign{}
\endhead
\bottomrule\noalign{}
\endlastfoot
\path{LAB3-IncidentResponse-HELIX.unl} & EVE-NG topology --- import directly \\
\href{https://github.com/TheEgea/TFG/tree/main/src/eve-ng/configs/nodes/Lab3/server-web}{\path{nodes/Lab3/server-web/}} & Pre-staged auth.log, bash history, sudoers \\
\href{https://github.com/TheEgea/TFG/tree/main/src/eve-ng/configs/nodes/Lab3/server-db}{\path{nodes/Lab3/server-db/}} & clients.db, exfil\_marker, network config \\
\href{https://github.com/TheEgea/TFG/tree/main/src/eve-ng/configs/nodes/Lab3/vyos}{\path{nodes/Lab3/vyos/}} & VyOS config.boot \\
\end{longtable}}

\section{LAB3 --- Infrastructure Reference}

Complete infrastructure reference for LAB3 HELIX Systems. For instructor/lab admin use.

\begin{center}\rule{0.5\linewidth}{0.5pt}\end{center}

\subsection{Node inventory}

{\footnotesize\begin{longtable}[]{@{}>{\raggedright\arraybackslash}p{0.17\linewidth}
  >{\raggedright\arraybackslash}p{0.17\linewidth}
  >{\raggedright\arraybackslash}p{0.17\linewidth}
  >{\raggedright\arraybackslash}p{0.17\linewidth}
  >{\raggedright\arraybackslash}p{0.17\linewidth}@{}}
\toprule\noalign{}
Node & EVE-NG ID & Console & Image & Status \\
\midrule\noalign{}
\endhead
\bottomrule\noalign{}
\endlastfoot
pfSense-LAB3 & 1 & VNC 32769 & pfsense-ce & Operational \\
VyOS-LAB3 & 2 & Telnet 32770 & vyos-rolling & Operational \\
Server-Web & 3 & VNC 32771 & linux-ubuntu-server-24.04 & Operational \\
Server-DB & 4 & VNC 32772 & linux-ubuntu-server-24.04 & Operational \\
\end{longtable}}

\subsection{Network addressing}

{\footnotesize\begin{longtable}[]{@{}>{\raggedright\arraybackslash}p{0.22\linewidth}
  >{\raggedright\arraybackslash}p{0.22\linewidth}
  >{\raggedright\arraybackslash}p{0.22\linewidth}
  >{\raggedright\arraybackslash}p{0.22\linewidth}@{}}
\toprule\noalign{}
Network & Subnet & Bridge & Host IP \\
\midrule\noalign{}
\endhead
\bottomrule\noalign{}
\endlastfoot
Net-Servers & 192.168.50.0/24 & vnet0\_2 & 192.168.50.254 \\
Net-Internal & 192.168.60.0/24 & vnet0\_3 & 192.168.60.254 \\
Net-Link (pfSense$\leftrightarrow$VyOS) & 172.16.2.0/30 & --- & --- \\
\end{longtable}}

\subsection{Credentials (full reference)}

{\footnotesize\begin{longtable}[]{@{}>{\raggedright\arraybackslash}p{0.22\linewidth}
  >{\raggedright\arraybackslash}p{0.22\linewidth}
  >{\raggedright\arraybackslash}p{0.22\linewidth}
  >{\raggedright\arraybackslash}p{0.22\linewidth}@{}}
\toprule\noalign{}
System & User & Password & Access \\
\midrule\noalign{}
\endhead
\bottomrule\noalign{}
\endlastfoot
EVE-NG host & root & SSH key id\_ed25519\_eve-ng & ssh eve-ng \\
pfSense-LAB3 & admin & pfsense & VNC 32769 \\
VyOS-LAB3 & vyos & vyos & Telnet 32770 \\
Server-Web & ubuntu & ubuntu & VNC 32771 (installer default) \\
Server-Web & analyst & An@lyst2024 & SSH via pfSense DNAT \\
Server-Web & devops & D3v0ps\#2023 & Compromised (attacker-used) \\
Server-Web & maint & --- & Backdoor account \\
Server-DB & ubuntu & ubuntu & VNC 32772 \\
Server-DB & analyst & An@lyst2024 & SSH from Server-Web \\
Server-DB & devops & D3v0ps\#2023 & Compromised (reused password) \\
\end{longtable}}

\subsection{Pre-staged evidence}

\subsubsection{\texorpdfstring{Server-Web ({\ttfamily 192.168.50.10})}{Server-Web (192.168.50.10)}}

{\footnotesize\begin{longtable}[]{@{}
  >{\raggedright\arraybackslash}p{0.17\linewidth}
  >{\raggedright\arraybackslash}p{0.54\linewidth}
  >{\raggedright\arraybackslash}p{0.19\linewidth}@{}}
\toprule\noalign{}
\begin{minipage}[b]{\linewidth}\raggedright
Artefact
\end{minipage} & \begin{minipage}[b]{\linewidth}\raggedright
Path
\end{minipage} & \begin{minipage}[b]{\linewidth}\raggedright
Content
\end{minipage} \\
\midrule\noalign{}
\endhead
\bottomrule\noalign{}
\endlastfoot
SSH brute-force + login & \path{/var/log/auth.log} & 25 failed + 1 successful login for \texttt{devops} from \path{185.220.101.47} \\
Privilege escalation & \path{/var/log/auth.log} & \path{sudo /usr/bin/find} by \texttt{devops} at 03:18 \\
Backdoor creation & \path{/var/log/auth.log} & \texttt{useradd\ maint} at 03:19 \\
Attacker command history & \path{/home/devops/.bash_history} & Full post-escalation command sequence \\
Backdoor account & \path{/etc/passwd} + \path{/etc/sudoers.d/maint} & \texttt{maint} user with \texttt{NOPASSWD:\ ALL} \\
\end{longtable}}

\subsubsection{\texorpdfstring{Server-DB ({\ttfamily 192.168.60.10})}{Server-DB (192.168.60.10)}}

{\footnotesize\begin{longtable}[]{@{}
  >{\raggedright\arraybackslash}p{0.30\linewidth}
  >{\raggedright\arraybackslash}p{0.30\linewidth}
  >{\raggedright\arraybackslash}p{0.30\linewidth}@{}}
\toprule\noalign{}
\begin{minipage}[b]{\linewidth}\raggedright
Artefact
\end{minipage} & \begin{minipage}[b]{\linewidth}\raggedright
Path
\end{minipage} & \begin{minipage}[b]{\linewidth}\raggedright
Content
\end{minipage} \\
\midrule\noalign{}
\endhead
\bottomrule\noalign{}
\endlastfoot
Lateral movement login & \path{/var/log/auth.log} & \texttt{devops} login from \path{192.168.50.10} at 03:21 \\
HELIX client database & \path{/opt/helix/data/clients.db} & SQLite DB with fictional client records \\
Exfiltration marker & \path{/opt/helix/data/.exfil_marker} & \texttt{185.220.101.47\ -\ 2026-04-24\ 03:22\ -\ data\ accessed} \\
\end{longtable}}

\subsection{Lab startup procedure}

\begin{lstlisting}
# 1. EVE-NG host --- restore bridge IPs after reboot
# (udev rule 99-lab3-bridges.rules should auto-apply)
# Verify:
ip addr show vnet0_2  # should have 192.168.50.254/24
ip addr show vnet0_3  # should have 192.168.60.254/24

# 2. Start all EVE-NG nodes (pfSense, VyOS, Server-Web, Server-DB)
# All are pre-configured --- just power on from EVE-NG UI

# 3. Verify student entry point
ssh analyst@<pfSense-WAN-IP>
# Should land on helix-web prompt
\end{lstlisting}

\begin{lstlisting}
The full student SSH chain (homelab $\rightarrow$ pfSense DNAT $\rightarrow$ Server-Web $\rightarrow$ Server-DB pivot) has been verified end-to-end in a live session.
\end{lstlisting}

\section{LAB3 --- VyOS Router}

VyOS-LAB3 provides inter-segment routing and SNAT masquerade for both internal networks. No DHCP or DNS forwarding --- all addresses are static.

\begin{center}\rule{0.5\linewidth}{0.5pt}\end{center}

\subsection{Interfaces}

{\footnotesize\begin{longtable}[]{@{}>{\raggedright\arraybackslash}p{0.22\linewidth}
  >{\raggedright\arraybackslash}p{0.22\linewidth}
  >{\raggedright\arraybackslash}p{0.22\linewidth}
  >{\raggedright\arraybackslash}p{0.22\linewidth}@{}}
\toprule\noalign{}
Interface & IP & Network & Description \\
\midrule\noalign{}
\endhead
\bottomrule\noalign{}
\endlastfoot
eth0 & 172.16.2.2/30 & Net-Link & Uplink to pfSense LAN \\
eth1 & 192.168.50.1/24 & Net-Servers & Gateway for Server-Web \\
eth2 & 192.168.60.1/24 & Net-Internal & Gateway for Server-DB \\
\end{longtable}}

\subsection{Access}

\begin{lstlisting}
# Console via EVE-NG telnet:
ssh eve-ng "telnet localhost 32770"
# user: vyos / vyos
\end{lstlisting}

\subsection{Interface and routing configuration}

\begin{lstlisting}
set interfaces ethernet eth0 address '172.16.2.2/30'
set interfaces ethernet eth0 description 'UPLINK-PFSENSE'
set interfaces ethernet eth1 address '192.168.50.1/24'
set interfaces ethernet eth1 description 'Net-Servers'
set interfaces ethernet eth2 address '192.168.60.1/24'
set interfaces ethernet eth2 description 'Net-Internal'

set protocols static route 0.0.0.0/0 next-hop 172.16.2.1
\end{lstlisting}

\subsection{NAT Source rules (SNAT masquerade)}

\begin{lstlisting}
set nat source rule 10 outbound-interface name 'eth0'
set nat source rule 10 source address '192.168.50.0/24'
set nat source rule 10 translation address 'masquerade'

set nat source rule 20 outbound-interface name 'eth0'
set nat source rule 20 source address '192.168.60.0/24'
set nat source rule 20 translation address 'masquerade'
\end{lstlisting}

\subsection{SSH service}

\begin{lstlisting}
set service ssh
commit
save
\end{lstlisting}

\begin{lstlisting}
Unlike LAB2, VyOS-LAB3 provides neither DHCP nor DNS forwarding. All node IPs are configured statically via Netplan.
\end{lstlisting}

\section{LAB3 --- pfSense Firewall}

pfSense-LAB3 is the perimeter firewall and student entry point. The key feature is a \textbf{DNAT rule} that forwards incoming TCP:22 directly to Server-Web, simulating the realistic scenario where an analyst connects to the affected host.

\begin{center}\rule{0.5\linewidth}{0.5pt}\end{center}

\subsection{Access}

{\footnotesize\begin{longtable}[]{@{}>{\raggedright\arraybackslash}p{0.30\linewidth}
  >{\raggedright\arraybackslash}p{0.30\linewidth}
  >{\raggedright\arraybackslash}p{0.30\linewidth}@{}}
\toprule\noalign{}
Interface & Assignment & IP \\
\midrule\noalign{}
\endhead
\bottomrule\noalign{}
\endlastfoot
vtnet0 & WAN & DHCP from homelab (\textasciitilde192.168.0.x) \\
vtnet1 & LAN & 172.16.2.1/30 (static) \\
\end{longtable}}

\subsection{Required fixes (applied at setup)}

Same two corrections as LAB2:

\begin{enumerate}
\def\labelenumi{\arabic{enumi}.}
\tightlist
\item
  \textbf{Block private networks} --- disabled on WAN (WAN shares 192.168.0.0/24 = RFC 1918 space)
\item
  \textbf{reply-to} directive --- removed from WAN rules (causes asymmetric routing in single-ISP setups)
\end{enumerate}

\subsection{DNAT rule --- student entry point}

Port forward on WAN interface:

{\footnotesize\begin{longtable}[]{@{}>{\raggedright\arraybackslash}p{0.45\linewidth}
  >{\raggedright\arraybackslash}p{0.45\linewidth}@{}}
\toprule\noalign{}
Field & Value \\
\midrule\noalign{}
\endhead
\bottomrule\noalign{}
\endlastfoot
Interface & WAN \\
Protocol & TCP \\
Destination port & 22 \\
Redirect target IP & 192.168.50.10 (Server-Web) \\
Redirect target port & 22 \\
\end{longtable}}

Students connect with:

\begin{lstlisting}
ssh analyst@<pfSense-WAN-IP>
# Password: An@lyst2024
\end{lstlisting}

The connection lands directly on \texttt{helix-web}, the primary forensic target.

\subsection{Static routes}

Two static routes pointing both internal segments via VyOS:

{\footnotesize\begin{longtable}[]{@{}>{\raggedright\arraybackslash}p{0.45\linewidth}
  >{\raggedright\arraybackslash}p{0.45\linewidth}@{}}
\toprule\noalign{}
Destination & Gateway \\
\midrule\noalign{}
\endhead
\bottomrule\noalign{}
\endlastfoot
192.168.50.0/24 & 172.16.2.2 (VyOS eth0) \\
192.168.60.0/24 & 172.16.2.2 (VyOS eth0) \\
\end{longtable}}

\begin{lstlisting}
pfSense CE outputs only to VGA framebuffer --- the EVE-NG telnet console (port 32769) is silent.
To interact: connect via raw socket and send `Ctrl-A C` to enter QEMU monitor mode.
Use `xp /40wx 0xb8000` to read the VGA text buffer and `sendkey` to inject keystrokes.
\end{lstlisting}

\section{LAB3 --- Server-Web (helix-web)}

Server-Web is the \textbf{primary forensic target}. It contains all the evidence of the initial attack phases: brute-force, successful login, privilege escalation, and backdoor creation.

\begin{center}\rule{0.5\linewidth}{0.5pt}\end{center}

\subsection{System}

{\footnotesize\begin{longtable}[]{@{}>{\raggedright\arraybackslash}p{0.45\linewidth}
  >{\raggedright\arraybackslash}p{0.45\linewidth}@{}}
\toprule\noalign{}
Parameter & Value \\
\midrule\noalign{}
\endhead
\bottomrule\noalign{}
\endlastfoot
Hostname & helix-web \\
OS & Ubuntu Server 24.04 LTS \\
IP & 192.168.50.10/24 \\
Gateway & 192.168.50.1 (VyOS eth1) \\
\end{longtable}}

\subsection{Network config (Netplan)}

\begin{lstlisting}
network:
  version: 2
  ethernets:
    ens3:
      addresses:
        - 192.168.50.10/24
      routes:
        - to: default
          via: 192.168.50.1
\end{lstlisting}

\subsection{OS accounts}

{\footnotesize\begin{longtable}[]{@{}
  >{\raggedright\arraybackslash}p{0.30\linewidth}
  >{\raggedright\arraybackslash}p{0.30\linewidth}
  >{\raggedright\arraybackslash}p{0.30\linewidth}@{}}
\toprule\noalign{}
\begin{minipage}[b]{\linewidth}\raggedright
Account
\end{minipage} & \begin{minipage}[b]{\linewidth}\raggedright
Password
\end{minipage} & \begin{minipage}[b]{\linewidth}\raggedright
Role
\end{minipage} \\
\midrule\noalign{}
\endhead
\bottomrule\noalign{}
\endlastfoot
ubuntu & ubuntu & Default installer account \\
analyst & An@lyst2024 & \textbf{Student entry} --- member of \texttt{adm} group (can read auth.log) \\
devops & D3v0ps\#2023 & Compromised developer account \\
maint & --- & \textbf{Backdoor} --- created by attacker at 03:19 \\
\end{longtable}}

\begin{lstlisting}
`/var/log/auth.log` is owned `root:adm` mode `640`. The `analyst` user must be in the `adm` group to read it. Applied via nbd chroot:
```bash
sudo chroot /mnt/lab3-web usermod -aG adm analyst
```
\end{lstlisting}

\subsection{Sudoers misconfiguration (privilege escalation vector)}

\begin{lstlisting}
# /etc/sudoers.d/devops
devops ALL=(ALL) NOPASSWD: /usr/bin/find
\end{lstlisting}

\texttt{find}'s \texttt{-exec} flag allows arbitrary command execution as root. The attacker ran:

\begin{lstlisting}
sudo /usr/bin/find / -name "*.conf" -exec cat  \;
\end{lstlisting}

\subsection{Pre-staged evidence in auth.log}

\begin{lstlisting}
# Brute-force phase (03:05--03:16): 25 entries like:
Apr 24 03:05:03 helix-web sshd[1201]: Failed password for devops from 185.220.101.47 port 42001 ssh2

# Successful login (03:17):
Apr 24 03:17:42 helix-web sshd[1231]: Accepted password for devops from 185.220.101.47 port 43100 ssh2

# Privilege escalation (03:18):
Apr 24 03:18:01 helix-web sudo: devops : TTY=pts/0 ; PWD=/home/devops ; USER=root ; COMMAND=/usr/bin/find

# Backdoor creation (03:19):
Apr 24 03:19:15 helix-web useradd[2041]: new user: name=maint, UID=1003, GID=1003, home=/home/maint, shell=/bin/bash
\end{lstlisting}

\subsection{Pre-staged command history (/home/devops/.bash\_history)}

\begin{lstlisting}
id
whoami
sudo /usr/bin/find / -name "*.conf" -exec cat  \;
useradd -m maint
echo "maint ALL=(ALL) NOPASSWD: ALL" > /etc/sudoers.d/maint
chmod 440 /etc/sudoers.d/maint
exit
\end{lstlisting}

\subsection{Useful investigation commands}

\begin{lstlisting}
# Count brute-force attempts
grep 'Failed password for devops' /var/log/auth.log | wc -l

# Find successful login + escalation
grep -E 'Accepted|sudo' /var/log/auth.log

# Find backdoor account creation
grep 'useradd\|maint' /var/log/auth.log

# Read attacker command history (needs sudo as analyst)
sudo cat /home/devops/.bash_history

# Check sudoers drop-in
sudo cat /etc/sudoers.d/devops
sudo cat /etc/sudoers.d/maint
\end{lstlisting}

\section{LAB3 --- Server-DB (helix-db)}

Server-DB is the \textbf{lateral movement and data exfiltration target}. It holds the HELIX client database and the exfiltration marker that confirms the attacker accessed sensitive data after pivoting from Server-Web.

\begin{center}\rule{0.5\linewidth}{0.5pt}\end{center}

\subsection{System}

{\footnotesize\begin{longtable}[]{@{}
  >{\raggedright\arraybackslash}p{0.45\linewidth}
  >{\raggedright\arraybackslash}p{0.45\linewidth}@{}}
\toprule\noalign{}
\begin{minipage}[b]{\linewidth}\raggedright
Parameter
\end{minipage} & \begin{minipage}[b]{\linewidth}\raggedright
Value
\end{minipage} \\
\midrule\noalign{}
\endhead
\bottomrule\noalign{}
\endlastfoot
Hostname & helix-db \\
OS & Ubuntu Server 24.04 LTS \\
IP & 192.168.60.10/24 \\
Gateway & 192.168.60.1 (VyOS eth2) \\
Segment & Net-Internal (isolated --- not directly reachable from homelab) \\
\end{longtable}}

\subsection{Network config (Netplan)}

\begin{lstlisting}
network:
  version: 2
  ethernets:
    ens3:
      addresses:
        - 192.168.60.10/24
      routes:
        - to: default
          via: 192.168.60.1
\end{lstlisting}

\subsection{Accounts}

{\footnotesize\begin{longtable}[]{@{}>{\raggedright\arraybackslash}p{0.30\linewidth}
  >{\raggedright\arraybackslash}p{0.30\linewidth}
  >{\raggedright\arraybackslash}p{0.30\linewidth}@{}}
\toprule\noalign{}
Account & Password & Role \\
\midrule\noalign{}
\endhead
\bottomrule\noalign{}
\endlastfoot
ubuntu & ubuntu & Default installer account \\
analyst & An@lyst2024 & Internal access (student pivot) \\
devops & D3v0ps\#2023 & \textbf{Reused password} --- same as Server-Web \\
\end{longtable}}

\subsection{Access from Server-Web}

\begin{lstlisting}
# From helix-web as analyst:
ssh analyst@192.168.60.10
# Password: An@lyst2024
\end{lstlisting}

\subsection{Pre-staged evidence in auth.log}

\begin{lstlisting}
# Lateral movement login (03:21):
Apr 24 03:21:03 helix-db sshd[987]: Accepted password for devops from 192.168.50.10 port 51234 ssh2
\end{lstlisting}

Source \path{192.168.50.10} = Server-Web internal IP --- confirms pivot from compromised host.

\subsection{HELIX data directory}

\begin{lstlisting}
ls -la /opt/helix/data/
\end{lstlisting}

\begin{lstlisting}
total 28
drwxr-xr-x 2 root root 4096 Apr 24 03:22 .
drwxr-xr-x 3 root root 4096 Apr 24 03:10 ..
-rw-r--r-- 1 root root 8192 Apr 24 03:10 clients.db
-rw-r--r-- 1 root root   47 Apr 24 03:22 .exfil_marker
\end{lstlisting}

\begin{lstlisting}
cat /opt/helix/data/.exfil_marker
# Output: 185.220.101.47 - 2026-04-24 03:22 - data accessed
\end{lstlisting}

The \path{.exfil_marker} timestamp (03:22) is consistent with the lateral movement login (03:21), confirming the attack sequence.

\texttt{clients.db} is an SQLite database with fictional HELIX Systems client records. Its presence establishes the data-access motive.

\subsection{Key finding}

The lateral movement succeeded purely due to \textbf{password reuse} --- \texttt{D3v0ps\#2023} was the same on both servers. Network segmentation (Server-DB is on Net-Internal, not directly reachable externally) did \textbf{not} prevent the attack once Server-Web was compromised.

\section{LAB3 --- Evidence Walkthrough}

Step-by-step investigation guide. Each step corresponds to specific log artefacts that the student is expected to locate, interpret, and document.

\begin{center}\rule{0.5\linewidth}{0.5pt}\end{center}

\subsection{Entry point}

\begin{lstlisting}
ssh analyst@<pfSense-WAN-IP>
# Password: An@lyst2024
# You land on: analyst@helix-web:~$
\end{lstlisting}

\begin{center}\rule{0.5\linewidth}{0.5pt}\end{center}

\subsection{Step 1 --- Establish context}

Before touching logs, orient yourself on the system.

\begin{lstlisting}
hostname
# helix-web

uname -a
# Linux helix-web 6.8.x-xx-generic

grep -v nologin /etc/passwd | grep -v false
# root, ubuntu, analyst, devops, maint
\end{lstlisting}

\begin{lstlisting}
The `maint` account is not a standard system account and has no documented business purpose. Note it immediately.
\end{lstlisting}

\begin{center}\rule{0.5\linewidth}{0.5pt}\end{center}

\subsection{Step 2 --- Brute-force identification}

\begin{lstlisting}
grep 'Failed password' /var/log/auth.log | head -5
# Apr 24 03:05:03 helix-web sshd[1201]: Failed password for devops from 185.220.101.47 ...
# Apr 24 03:05:06 helix-web sshd[1202]: Failed password for devops from 185.220.101.47 ...

grep 'Failed password for devops' /var/log/auth.log | wc -l
# 25
\end{lstlisting}

25 consecutive failures from \path{185.220.101.47} across a sustained window = \textbf{brute-force attack}.

\begin{center}\rule{0.5\linewidth}{0.5pt}\end{center}

\subsection{Step 3 --- Successful login and privilege escalation}

\begin{lstlisting}
grep -E 'Accepted|sudo' /var/log/auth.log
# Apr 24 03:17:42 helix-web sshd[1231]: Accepted password for devops from 185.220.101.47 ...
# Apr 24 03:18:01 helix-web sudo: devops : ... COMMAND=/usr/bin/find
\end{lstlisting}

The gap between the last failed attempt and the successful login suggests the attacker switched to the correct credential shortly after the brute-force phase ended. \path{sudo /usr/bin/find} by a developer account is anomalous --- \texttt{find} has no legitimate need for root.

\begin{center}\rule{0.5\linewidth}{0.5pt}\end{center}

\subsection{Step 4 --- Privilege escalation vector}

\begin{lstlisting}
sudo cat /etc/sudoers.d/devops
# devops ALL=(ALL) NOPASSWD: /usr/bin/find
\end{lstlisting}

\texttt{find\ -exec} allows running arbitrary commands as root. This is a well-known GTFOBins vector.

\begin{center}\rule{0.5\linewidth}{0.5pt}\end{center}

\subsection{Step 5 --- Backdoor account and persistence}

\begin{lstlisting}
grep 'useradd\|maint' /var/log/auth.log
# Apr 24 03:19:15 helix-web useradd[2041]: new user: name=maint, UID=1003 ...

sudo cat /home/devops/.bash_history
# id
# whoami
# sudo /usr/bin/find / -name "*.conf" -exec cat  \;
# useradd -m maint
# echo "maint ALL=(ALL) NOPASSWD: ALL" > /etc/sudoers.d/maint
# chmod 440 /etc/sudoers.d/maint
# exit

sudo cat /etc/sudoers.d/maint
# maint ALL=(ALL) NOPASSWD: ALL
\end{lstlisting}

The \path{.bash_history} confirms the exact sequence. The attacker created a backdoor account with full unrestricted sudo.

\begin{center}\rule{0.5\linewidth}{0.5pt}\end{center}

\subsection{Step 6 --- Lateral movement to Server-DB}

Pivot to Server-DB:

\begin{lstlisting}
ssh analyst@192.168.60.10
# Password: An@lyst2024
# analyst@helix-db:~$
\end{lstlisting}

Check the auth log:

\begin{lstlisting}
grep 'Accepted' /var/log/auth.log
# Apr 24 03:21:03 helix-db sshd[987]: Accepted password for devops from 192.168.50.10 port 51234 ssh2
\end{lstlisting}

Source \path{192.168.50.10} = Server-Web. The attacker reused \texttt{devops} credentials to reach the database server.

\begin{lstlisting}
You connect as `analyst` (your investigator account). The attacker previously moved laterally as `devops` --- hence the `auth.log` entry shows `devops`, not `analyst`.
\end{lstlisting}

\begin{center}\rule{0.5\linewidth}{0.5pt}\end{center}

\subsection{Step 7 --- Data exfiltration}

\begin{lstlisting}
ls -la /opt/helix/data/
# clients.db  .exfil_marker

cat /opt/helix/data/.exfil_marker
# 185.220.101.47 - 2026-04-24 03:22 - data accessed
\end{lstlisting}

The external attacker IP (\path{185.220.101.47}) matches the brute-force source. Timestamp \texttt{03:22} follows the lateral movement login at \texttt{03:21}.

\begin{center}\rule{0.5\linewidth}{0.5pt}\end{center}

\subsection{Incident findings summary}

{\footnotesize\begin{longtable}[]{@{}
  >{\raggedright\arraybackslash}p{0.20\linewidth}
  >{\raggedright\arraybackslash}p{0.20\linewidth}
  >{\raggedright\arraybackslash}p{0.28\linewidth}
  >{\raggedright\arraybackslash}p{0.22\linewidth}@{}}
\toprule\noalign{}
\begin{minipage}[b]{\linewidth}\raggedright
Time
\end{minipage} & \begin{minipage}[b]{\linewidth}\raggedright
Host
\end{minipage} & \begin{minipage}[b]{\linewidth}\raggedright
Artefact
\end{minipage} & \begin{minipage}[b]{\linewidth}\raggedright
Finding
\end{minipage} \\
\midrule\noalign{}
\endhead
\bottomrule\noalign{}
\endlastfoot
03:05--03:16 & Server-Web & auth.log & 25 failed SSH attempts (brute force) from \path{185.220.101.47} \\
03:17 & Server-Web & auth.log & Successful login as \texttt{devops} \\
03:18 & Server-Web & auth.log & Privilege escalation via \texttt{sudo\ find} \\
03:19 & Server-Web & auth.log + /etc/passwd & Backdoor account \texttt{maint} created \\
03:21 & Server-DB & auth.log & Lateral movement from \path{192.168.50.10} \\
03:22 & Server-DB & .exfil\_marker & Access to \texttt{clients.db} confirmed \\
\end{longtable}}

\begin{center}\rule{0.5\linewidth}{0.5pt}\end{center}

\subsection{Mitigations}

{\footnotesize\begin{longtable}[]{@{}
  >{\raggedright\arraybackslash}p{0.30\linewidth}
  >{\raggedright\arraybackslash}p{0.30\linewidth}
  >{\raggedright\arraybackslash}p{0.30\linewidth}@{}}
\toprule\noalign{}
\begin{minipage}[b]{\linewidth}\raggedright
Attack phase
\end{minipage} & \begin{minipage}[b]{\linewidth}\raggedright
Root cause
\end{minipage} & \begin{minipage}[b]{\linewidth}\raggedright
Fix
\end{minipage} \\
\midrule\noalign{}
\endhead
\bottomrule\noalign{}
\endlastfoot
Brute force & No rate limiting & fail2ban; SSH key auth only; disable password auth \\
Initial access & Weak/reused password & Strong password policy; MFA \\
Privilege escalation & \path{NOPASSWD: /usr/bin/find} & Never grant NOPASSWD to utility binaries \\
Persistence & Unrestricted root access & Audit \path{/etc/passwd} and \path{sudoers.d/} with FIM \\
Lateral movement & Password reuse & Unique credentials per system \\
Data exfiltration & No access control on \path{/opt/helix/data/} & Least-privilege ownership; audit logging \\
\end{longtable}}
